\documentclass[10pt,a4paper]{amsart}
\usepackage[margin=19mm]{geometry}
\usepackage{amsmath,amssymb,mathtools}
\usepackage[scr=rsfs,cal=euler]{mathalfa}
\usepackage{xfrac}
\usepackage[unicode,hidelinks,bookmarks=false]{hyperref}
\newcommand{\ie}{i.e.\ }
\newcommand{\cf}{cf.\ }
\newcommand{\resp}{respectively\ }
\newcommand{\Subsection}{\subsection}
\providecommand{\nobreakdash}{\nobreak\hbox{-}}
\setlength{\emergencystretch}{2em}
\multlinegap0pt
\usepackage{mathtools}
\usepackage{amssymb}
\usepackage{algorithm}\usepackage{algpseudocode}
\usepackage{tikz}
\usepackage[normalem]{ulem}
\usepackage[shortlabels]{enumitem}
\newtheorem{corollary}{Corollary}
\newtheorem{lemma}{Lemma}
\usepackage{cleveref}
\crefname{corollary}{Corollary}{Corollarys}
\crefname{lemma}{Lemma}{Lemmas}
\newcommand\R{{\mathbb R}}
\newcommand{\bA}{{\bf A}}
\newcommand{\bI}{{\bf Id}}
\newcommand{\norm}[1]{\left\Vert#1\right\Vert}
\title{Approximately Dual and Pseudo-Dual Probabilistic Frames}
\author{Dongwei Chen, Emily J. King, Clayton Shonkwiler}
\date{Source: arXiv:2505.13885v2 (CC BY 4.0)}
\begin{document}
\maketitle

\noindent\emph{Source and licence.} Approximately Dual and Pseudo-Dual Probabilistic Frames. Version arXiv:2505.13885v2 (CC BY 4.0), distributed by arXiv under the Creative Commons Attribution 4.0 International licence, http://creativecommons.org/licenses/by/4.0/, as recorded in the arXiv licence metadata for this exact version and shown on the abstract page https://arxiv.org/abs/2505.13885v2. Full licence notice: \texttt{licenses/arxiv-2505.13885-CC-BY-4.0.txt}.\par\smallskip

\noindent\textit{Editorial scope.} Counted unit: the article's lemma DualIsFrame (source0.tex lines 440--453). The article's complete original proof of it is delivered with it (source0.tex lines 455--545, one \texttt{proof} environment of 91 lines). No mathematics is written, repaired or re-derived: the statement and the proof are the article's own lines, in the article's own notation, and no retained line is substituted or re-flowed.  Retained uncounted as the source-local context the proof needs: lines 203--205.  Nothing was omitted from any retained span, so the package carries no omission record.  No retained line contains a citation, so the wrapper carries no bibliography and no bibliography entry.  No retained line contains a figure environment, an image-inclusion command or drawing code, so no image is delivered and the package declares no asset dependency.  Editorial formatting only, in addition: an AMS \texttt{amsart} preamble that restates the article's own declarations and macros used by the retained lines, namely the environments \texttt{corollary}, \texttt{lemma} on separate counters, together with the standard packages those lines need.  The retained environments print in this document as Corollary 1, Lemma 1 in order of appearance, whereas the article prints the same items with the numbering of its own full document.  The complete native source of the article as distributed in the arXiv e-print for this exact version is bundled unchanged as \texttt{sources/178404/source0.tex}.
\begin{corollary}\label{cor:full measure contains basis}
    If $\mu$ is a probabilistic frame for $\R^n$, then any Borel set $E \subseteq \R^n$ with $\mu(E) = 1$ contains a basis for $\R^n$.
\end{corollary}

\begin{lemma}\label{DualIsFrame}
Let $\mu$ be a probabilistic frame with upper frame bound $B>0$.  Suppose $\nu \in \mathcal{P}_2(\mathbb{R}^n)$ is an approximate dual of $\mu$ with mixed frame operator $\bf A$, then $\nu$ is a probabilistic frame with bounds $ 1/(B \| {\bf A}^{-1}\|^2)$ and $M_2(\nu)$. 

In addition, for any ${\bf f} \in \mathbb{R}^n$, we have the following uncertainty principle:
     \begin{equation*}
        \int_{\mathbb{R}^n} |\langle {\bf A }^{-1} {\bf x}, {\bf f} \rangle |^2 d\mu({\bf x}) \int_{\mathbb{R}^n} |\langle {\bf y}, {\bf f} \rangle |^2 d\nu({\bf y})  \geq \| {\bf f} \|^4,
    \end{equation*}
and equality is attained for all ${\bf f} \in \mathbb{R}^n$  if and only if  ${\bf S}_{\mu}=c{\bf A}{\bf A}^{t}$ and
$\nu=(\frac{1}{c}{\bf A}^{-1})_{\#}\mu$ where $c$ is a nonzero constant.
Furthermore, if $A_\mu$, $B_\mu$, and $A_\nu$, $B_\nu$ are optimal frame bounds for $\mu$ and $\nu$, then 
\begin{equation*}
    A_\nu \geq \frac{1}{B_\mu \| {\bf A}^{-1}\|^2} \ \text{and} \  A_\mu \geq \frac{1}{B_\nu \| {\bf A}^{-1}\|^2}.
\end{equation*}
\end{lemma}

\begin{proof}
 Since $\nu \in \mathcal{P}_2(\mathbb{R}^n)$ is an approximate dual of $\mu$ with mixed frame operator $\bf A$, then  $\bf A$ is invertible with $\norm{\bA - \bI}<1$, and there exists $\gamma \in \Gamma(\mu, \nu)$ such that 
   $$   \int_{\mathbb{R}^n \times \mathbb{R}^n} {\bf x} {\bf y}^t d\gamma({\bf x, y}) = {\bf A}.$$
   Then, for any ${\bf f} \in \mathbb{R}^n$, 
   \begin{equation*}
   \begin{split}
       \|{\bf f}\|^4 = |\langle {\bf f}, {\bf f} \rangle|^2 = |\langle {\bf A}^{-1}{\bf A}{\bf f}, {\bf f} \rangle|^2 &= \left|\int_{\mathbb{R}^n \times \mathbb{R}^n} \langle {\bf A}^{-1} {\bf x}, {\bf f} \rangle \langle {\bf y,f} \rangle d\gamma({\bf x, y}) \right |^2 \\
       & \leq \int_{\mathbb{R}^n} |\langle {\bf A }^{-1} {\bf x}, {\bf f} \rangle |^2 d\mu({\bf x}) \int_{\mathbb{R}^n} |\langle {\bf y}, {\bf f} \rangle |^2 d\nu({\bf y}) \\
       & \leq B \|({\bf A }^{-1})^t {\bf f}\|^2 \int_{\mathbb{R}^n} |\langle {\bf y}, {\bf f} \rangle |^2 d\nu({\bf y}) \\
       &\leq B\|{\bf A }^{-1}\|^2 \| {\bf f}\|^2 \int_{\mathbb{R}^n} |\langle {\bf y}, {\bf f} \rangle |^2 d\nu({\bf y}),
   \end{split}
   \end{equation*}
where the first inequality is due to the Cauchy–Schwarz inequality and the second follows from the frame property of $\mu$. Note that we have proved the desired uncertainty principle along the way. Therefore, for any ${\bf f} \in \mathbb{R}^n$,
\begin{equation*}
       \frac{\|{\bf f}\|^2}{B \|{\bf A }^{-1}\|^2}
       \leq  \int_{\mathbb{R}^n} |\langle {\bf y}, {\bf f} \rangle |^2 d\nu({\bf y}) \leq M_2(\nu)\|{\bf f}\|^2,
   \end{equation*}
 which implies that $\nu$ is a probabilistic frame with bounds $ 1/(B \| {\bf A}^{-1}\|^2)$ and $M_2(\nu)$. 

In addition, for a fixed ${\bf f}$, equality holds in the uncertainty principle if and only if the two scalar functions
\[
({\bf x},{\bf y})\longmapsto \langle {\bf A}^{-1}{\bf x},{\bf f}\rangle
\quad\text{and}\quad
({\bf x},{\bf y})\longmapsto \langle {\bf y},{\bf f}\rangle
\]
are linearly dependent in $L^{2}(\gamma, \mathbb{R}^n \times \mathbb{R}^n)$; that is, there exists a scalar
$c({\bf f})$ that may depend on ${\bf f}$ such that
\begin{equation}\label{eq:inner product ae}
\langle {\bf A}^{-1}{\bf x},{\bf f}\rangle
= c({\bf f}) \langle {\bf y},{\bf f}\rangle
\ \text{for $\gamma$-almost every }({\bf x},{\bf y}).
\end{equation}

We claim that if the equality is true for all $\mathbf{f} \in \mathbb{R}^n$, then $c({\bf f})$ is a constant, independent of ${\bf f}$. 

To see this, note first that $c(\alpha {\bf f}) = c({\bf f})$ for all $\alpha \neq 0$. 

Now, let $E$ be the set on which~\eqref{eq:inner product ae} holds. Then $E$ has full $\gamma$-measure and hence the projection $\pi_y(E)$ has full $\nu$-measure. Therefore, by \cref{cor:full measure contains basis} there exist ${\bf x}_1, \dots , {\bf x}_n, {\bf y}_1, \dots {\bf y}_n \in \R^n$ so that ${\bf y}_1, \dots {\bf y}_n$ is a basis for $\R^n$ and $({\bf x}_i, {\bf y}_i) \in E$ for all $i=1, \dots , n$.

For any $i=1, \dots , n$, suppose ${\bf f} \in \R^n$ is a unit vector with $\langle {\bf y}_i, {\bf f} \rangle \neq 0$. Write 
\[
    {\bf f} = \alpha {\bf y}_i + {\bf g},
\]
where $\alpha \neq 0$ and ${\bf g} \in {\bf y}_i^\perp$. Then
\[
\langle {\bf A}^{-1}{\bf x}_i,\alpha {\bf y}_i\rangle
= c(\alpha {\bf y}_i) \langle {\bf y}_i,\alpha {\bf y}_i\rangle = c({\bf y}_i) \langle {\bf y}_i,\alpha {\bf y}_i\rangle
\]
and
\[
    \langle {\bf A}^{-1}{\bf x}_i,{\bf g}\rangle = c({\bf g}) \langle {\bf y}_i,{\bf g}\rangle = 0
\]
Therefore,
\[
    c({\bf f}) \alpha \|{\bf y}_i\|^2 = c({\bf f})\langle {\bf y}_i, {\bf f}\rangle = \langle {\bf A}^{-1} {\bf x}_i , {\bf f} \rangle = \langle {\bf y}_i , c({\bf y}_i) \alpha {\bf y}_i\rangle = c({\bf y}_i) \alpha \|{\bf y}_i\|^2,
\]
so $c({\bf f}) = c({\bf y}_i)$. 

In other words, $c$ is constant on unit vectors in $\R^n \backslash{\bf y}_i^\perp$ and hence, by the scale-invariance of $c$, on all of each $\R^n \backslash{\bf y}_i^\perp$. Since these sets have nontrivial intersections, we see that all constants are the same, so $c$ is constant on 
$$\bigcup_{i=1}^n \left(\R^n \backslash {\bf y}_i^\perp\right) = \R^n \backslash \left(\bigcap_{i=1}^n {\bf y}_i^\perp\right).
$$
Since ${\bf y}_1, \dots {\bf y}_n$ is a basis for $\R^n$, the intersection $\bigcap_{i=1}^n {\bf y}_i^\perp = \{{\bf 0}\}$, so $c$ is constant on $\R^n \backslash \{{\bf 0}\}$.

Moreover, re-defining $c({\bf 0})$ to equal this common constant doesn't affect the validity of~\eqref{eq:inner product ae}, so we see that $c$ can be chosen to be constant on all of $\R^n$.

Therefore, we have shown if equality in the uncertainty principle is true for all ${\bf f} \in \mathbb{R}^n$,  then
$$\langle {\bf A}^{-1}{\bf x},{\bf f}\rangle
= c \langle {\bf y},{\bf f}\rangle
\ \text{for $\gamma$-almost every }({\bf x},{\bf y}),$$ which implies that $\gamma = (\mathbf{Id}, \frac{1}{c}{\bf A}^{-1})_\#\mu$. Since $\gamma \in \Gamma(\mu, \nu)$, then $\nu = (\frac{1}{c}{\bf A}^{-1})_\#\mu$. Also $\int_{\mathbb{R}^n \times \mathbb{R}^n} {\bf x} {\bf y}^t d\gamma({\bf x, y}) = {\bf A}$ implies that ${\bf S}_{\mu}=c{\bf A}{\bf A}^{t}$.

On the other hand, suppose ${\bf S}_{\mu}=c{\bf A}{\bf A}^{t}$ and
$\nu=(\frac{1}{c}{\bf A}^{-1})_{\#}\mu$ for some nonzero constant $c$.
Then $\nu$ is an approximately dual frame of $\mu$ with mixed frame operator ${\bf A}$ against $\gamma = (\mathbf{Id}, \frac{1}{c}{\bf A}^{-1})_\#\mu \in \Gamma(\mu, \nu)$, which follows from 
   $$   \int_{\mathbb{R}^n \times \mathbb{R}^n} {\bf x} {\bf y}^t d\gamma({\bf x, y}) = \frac{1}{c} \int_{\mathbb{R}^n} {\bf x} {\bf x}^t d\mu({\bf x}) ({\bf A}^{-1})^t = \frac{1}{c} {\bf S}_{\mu} ({\bf A}^{-1})^t = {\bf A}.
   $$
Since  $\gamma = (\mathbf{Id}, \frac{1}{c}{\bf A}^{-1})_\#\mu$, then for $\gamma$-almost every $({\bf x},{\bf y})$, ${\bf y}=\frac{1}{c}{\bf A}^{-1}{\bf x}$. Thus for any fixed ${\bf f}$, 
$\langle {\bf A}^{-1}{\bf x},{\bf f}\rangle
= c \langle {\bf y},{\bf f}\rangle$, for $\gamma$-almost every $({\bf x},{\bf y})$. 
Therefore, equality in the uncertainty principle holds for all $\mathbf{f} \in \mathbb{R}^n$. 


 Finally, if $A_\mu$, $B_\mu$ and $A_\nu$, $B_\nu$ are optimal bounds for $\mu$ and $\nu$, then the above result gives $ A_\nu \geq \frac{1}{B_\mu \| {\bf A}^{-1}\|^2}$, and similarly,  for any ${\bf f} \in \mathbb{R}^n$, we have 
   \begin{equation*}
   \begin{split}
       \|{\bf f}\|^4  = |\langle ({\bf A}^{-1})^t{\bf A}^t{\bf f}, {\bf f} \rangle|^2 &= \left|\int_{\mathbb{R}^n \times \mathbb{R}^n} \langle ({\bf A}^{-1})^t {\bf y}, {\bf f} \rangle \langle {\bf x,f} \rangle d\gamma({\bf x, y}) \right |^2 \\
       & \leq  \int_{\mathbb{R}^n} |\langle {\bf y}, {\bf A}^{-1} {\bf f} \rangle |^2 d\nu({\bf y}) \int_{\mathbb{R}^n} |\langle {\bf x}, {\bf f} \rangle |^2 d\mu({\bf x})\\
       &\leq B_\nu \|{\bf A }^{-1}\|^2 \| {\bf f}\|^2 \int_{\mathbb{R}^n} |\langle {\bf x}, {\bf f} \rangle |^2 d\mu({\bf x}).
   \end{split}
   \end{equation*}
Therefore, $\frac{1}{B_\nu \| {\bf A}^{-1}\|^2}$ is a lower frame bound for $\mu$ and thus $A_\mu \geq \frac{1}{B_\nu \| {\bf A}^{-1}\|^2}$.
\end{proof}

\end{document}
